\documentclass[12pt]{article}
\usepackage[ukrainian]{babel}
\usepackage{fontspec}
% --- НАЛАШТУВАННЯ ШРИФТІВ ---
\setmainfont{Schoolbook-Regular.otf}[
    Path = ../../,
    BoldFont = Schoolbook-Bold.otf,
    ItalicFont = Schoolbook-Italic.otf,
    BoldItalicFont = Schoolbook-BoldItalic.otf
]
\usepackage[italic]{mathastext}
\usepackage[a4paper,margin=1.8cm,bottom=2cm,top=2cm]{geometry}
\usepackage{amsmath,amssymb}
\usepackage{tikz}
\usepackage{pgfplots}
\pgfplotsset{compat=1.16}
\usetikzlibrary{calc,patterns,angles,quotes}
\usepackage{xcolor,array,fancyhdr,enumitem,multicol}
\usepackage{colortbl}
\usepackage{tcolorbox}
\tcbuselibrary{skins}
\usepackage[hidelinks]{hyperref}
\usepackage[firstpage=false, text=@pvtr2525, scale=0.6, color=gray!12]{draftwatermark}
\definecolor{mainGreen}{RGB}{34, 120, 64}
\definecolor{yearOrange}{RGB}{220, 100, 30}
\definecolor{yearcolor}{RGB}{220, 100, 30}
\definecolor{titleBg}{RGB}{225, 240, 220}
\definecolor{instrBg}{RGB}{255, 240, 220}
\definecolor{instrBorder}{RGB}{220, 150, 80}
\pagestyle{fancy}
\fancyhf{}
\renewcommand{\headrulewidth}{0.4pt}
\renewcommand{\footrulewidth}{0.4pt}
\fancyhead[C]{\small\color{gray!80} Збірник НМТ \textendash{} сесії 23.05--5.06.2026}
\fancyhead[R]{\small\color{mainGreen}\textbf{\href{https://t.me/pvtr2525}{@pvtr2525}}}
\fancyfoot[L]{\small\color{gray!80}\href{https://t.me/pvtr2525}{ТГ: @pvtr2525}}
\fancyfoot[C]{\small\color{gray!80}\thepage}
\fancyfoot[R]{\small\color{gray!80}НМТ 2026}
% --- ТАБЛИЦІ ВІДПОВІДЕЙ ---
\newcommand{\answerTable}[5]{%
\par\vspace{0.15cm}
\noindent
\begin{tabular}{|*{5}{>{\centering\arraybackslash}m{3cm}|}}
\hline
\rule[-0.15cm]{0pt}{0.6cm}\textbf{А} & \rule[-0.15cm]{0pt}{0.6cm}\textbf{Б} & \rule[-0.15cm]{0pt}{0.6cm}\textbf{В} & \rule[-0.15cm]{0pt}{0.6cm}\textbf{Г} & \rule[-0.15cm]{0pt}{0.6cm}\textbf{Д} \\
\hline
\rule[-0.2cm]{0pt}{0.75cm}#1 & \rule[-0.2cm]{0pt}{0.75cm}#2 & \rule[-0.2cm]{0pt}{0.75cm}#3 & \rule[-0.2cm]{0pt}{0.75cm}#4 & \rule[-0.2cm]{0pt}{0.75cm}#5 \\
\hline
\end{tabular}
\par\vspace{0.25cm}
}
\newcommand{\answerTableTall}[5]{%
\par\vspace{0.15cm}
\noindent
\begin{tabular}{|*{5}{>{\centering\arraybackslash}m{3cm}|}}
\hline
\rule[-0.2cm]{0pt}{0.7cm}\textbf{А} & \rule[-0.2cm]{0pt}{0.7cm}\textbf{Б} & \rule[-0.2cm]{0pt}{0.7cm}\textbf{В} & \rule[-0.2cm]{0pt}{0.7cm}\textbf{Г} & \rule[-0.2cm]{0pt}{0.7cm}\textbf{Д} \\
\hline
\rule[-0.45cm]{0pt}{1.2cm}#1 & \rule[-0.45cm]{0pt}{1.2cm}#2 & \rule[-0.45cm]{0pt}{1.2cm}#3 & \rule[-0.45cm]{0pt}{1.2cm}#4 & \rule[-0.45cm]{0pt}{1.2cm}#5 \\
\hline
\end{tabular}
\par\vspace{0.25cm}
}
\newcommand{\instructionBox}[1]{%
\par\vspace{0.3cm}
\begin{tcolorbox}[colback=instrBg, colframe=instrBorder, boxrule=0.6pt, arc=2pt,
    left=8pt, right=8pt, top=4pt, bottom=4pt]
\centering\small\bfseries #1
\end{tcolorbox}
\par\vspace{0.15cm}
}
\newcommand{\sectionTitle}[1]{%
\par\vspace{0.4cm}
\begin{tcolorbox}[colback=titleBg, colframe=titleBg, boxrule=0pt, arc=8pt,
    halign=center, fontupper=\Large\bfseries\color{mainGreen}]
#1
\end{tcolorbox}
\par\vspace{0.3cm}
}
\newcommand{\task}[2]{\noindent\textbf{#1.}\ \ #2\par\vspace{0.2cm}}
% --- НМТ-СТИЛЬ ПОЛЕ ВІДПОВІДІ ---
\newcommand{\nmtAnswerBox}{%
\par\vspace{0.2cm}\noindent
Відповідь:\ \
\begin{tabular}{@{}*{4}{|>{\centering\arraybackslash}p{0.8cm}}|@{\hspace{0.1cm}\textbf{,}\hspace{0.1cm}}*{3}{|>{\centering\arraybackslash}p{0.8cm}}|@{}}
\hline
\rule[-0.2cm]{0pt}{0.7cm} & & & & & & \\
\hline
\end{tabular}
\par\vspace{0.3cm}
}
% --- СІТКА ДЛЯ МАТЧИНГУ ---
\newcommand{\matchingGrid}{%
    \begingroup
    \renewcommand{\arraystretch}{1.15}
    \setlength{\tabcolsep}{4pt}
    \footnotesize
    \begin{tabular}{r|c|c|c|c|c|}
         \multicolumn{1}{c}{} & \multicolumn{1}{c}{\textbf{А}} & \multicolumn{1}{c}{\textbf{Б}} & \multicolumn{1}{c}{\textbf{В}} & \multicolumn{1}{c}{\textbf{Г}} & \multicolumn{1}{c}{\textbf{Д}} \\ \cline{2-6}
         \textbf{1} & \phantom{X} & \phantom{X} & \phantom{X} & \phantom{X} & \phantom{X} \\ \cline{2-6}
         \textbf{2} & & & & & \\ \cline{2-6}
         \textbf{3} & & & & & \\ \cline{2-6}
    \end{tabular}
    \endgroup
}
% --- ДОДАТКОВІ МАКРОСИ ЗБІРНИКА ---
\newcommand{\taskBlock}[1]{%
\par\vspace{0.45cm}
\noindent{\color{mainGreen}\rule{\linewidth}{0.8pt}}\par\vspace{0.12cm}
\noindent{\small\bfseries\color{mainGreen}#1}\par\vspace{0.25cm}
}
\newcounter{zad}
\newcommand{\zadtask}[1]{\stepcounter{zad}\noindent\textbf{\thezad.}\ \ #1\par\vspace{0.2cm}}
\newcommand{\zadnum}{\stepcounter{zad}\textbf{\thezad.}\ \ }
\newcounter{ans}
\newcommand{\ansitem}[1]{\stepcounter{ans}\noindent\textbf{\theans.}\ #1\par\vspace{0.07cm}}
\newcommand{\nmtyear}[1]{\hfill{\small\color{yearOrange}(НМТ #1)}}% --- ДОДАТКОВІ МАКРОСИ (розв'язки, зміст) ---
\tcbuselibrary{breakable}
\newcommand{\solution}[1]{%
\par\vspace{0.12cm}
\begin{tcolorbox}[colback=titleBg!45, colframe=mainGreen!55, boxrule=0.4pt, arc=2pt,
    left=6pt, right=6pt, top=3pt, bottom=3pt, breakable]
\footnotesize\textbf{Короткий розв'язок (оригінал).}\ #1
\end{tcolorbox}
\par\vspace{0.12cm}
}
\setcounter{tocdepth}{2}
\newcommand{\chapterTitle}[1]{%
\clearpage\phantomsection\addcontentsline{toc}{section}{#1}%
\sectionTitle{#1}}
\newcommand{\typeTitle}[1]{%
\phantomsection\addcontentsline{toc}{subsection}{\quad #1}%
\par\vspace{0.3cm}
\noindent{\large\bfseries\color{yearOrange}#1}\par\vspace{0.08cm}
\noindent{\color{yearOrange}\rule{\linewidth}{0.4pt}}\par\vspace{0.2cm}}
\newcommand{\ansTheme}[1]{\par\vspace{0.25cm}\noindent{\bfseries\color{mainGreen}#1}\par\vspace{0.12cm}}
\newcommand{\ansType}[1]{\par\vspace{0.1cm}\noindent{\itshape\color{yearOrange}#1}\par\vspace{0.08cm}}
\newcommand{\solbtn}[1]{\hypertarget{task:#1}{}\par\vspace{0.05cm}\noindent\hyperlink{sol:#1}{\fbox{\footnotesize пояснення $\rightarrow$}}\par\vspace{0.2cm}}
\begin{document}
\thispagestyle{empty}
\vspace*{1.2cm}
\begin{center}
{\Large\bfseries\color{mainGreen} ОРИГІНАЛЬНІ ЗАВДАННЯ НМТ-2026}\\[0.4cm]
{\Huge\bfseries\color{mainGreen} РОЗДІЛ 2. РІВНЯННЯ І НЕРІВНОСТІ}\\[0.6cm]
{\large Лише оригінали сесій 23.05--5.06.2026, без аналогів}\\[1cm]
{\large\color{mainGreen}\textbf{\href{https://t.me/pvtr2525}{Телеграм: @pvtr2525}}}
\end{center}
\clearpage
\chapterTitle{РОЗДІЛ 2. РІВНЯННЯ І НЕРІВНОСТІ}
\typeTitle{Лінійні рівняння та нерівності}
\taskBlock{Найпростіше лінійне рівняння --- сесія 23.05, завдання №1}
\zadtask{Розв'яжіть рівняння $\dfrac{x}{7} = 9$. \nmtyear{2026}}
\answerTableTall{$2$}{$63$}{$-63$}{$16$}{$\dfrac{9}{7}$}
\solbtn{1}
\taskBlock{Текстова задача на складання рівняння --- сесія 23.05, завдання №5}
\zadtask{У коробці $30$ тістечок двох видів: бісквітні й безе. Скільки всього бісквітних тістечок, якщо їх у $5$ разів більше, ніж безе? \nmtyear{2026}}
\answerTable{$24$}{$25$}{$6$}{$15$}{$5$}
\solbtn{2}
\taskBlock{Подвійна нерівність --- сесія 2.06, завдання №1}
\zadtask{Розв'яжіть нерівність $-8 \leqslant x - 3 \leqslant 5$.}
\answerTable{$[-11;\, 2]$}{$[-5;\, 8]$}{$[-5;\, 2]$}{$[-8;\, 5]$}{$[-11;\, 8]$}
\solbtn{3}
\taskBlock{Лінійна нерівність --- сесія 3.06, завдання №1}
\zadtask{Розв'яжіть нерівність $2x - 1 \geqslant 0$.}
\answerTableTall{$\left[-\dfrac{1}{2};\, +\infty\right)$}{$\left(-\infty;\, \dfrac{1}{2}\right]$}{$[1;\, +\infty)$}{$\left(-\infty;\, -\dfrac{1}{2}\right]$}{$\left[\dfrac{1}{2};\, +\infty\right)$}
\solbtn{4}
\typeTitle{Системи рівнянь}
\taskBlock{Система рівнянь --- сесія 30.05, завдання №1}
\zadtask{Розв'яжіть систему рівнянь $\begin{cases} \dfrac{2}{x} = -0{,}5 \\[0.2cm] \dfrac{y}{3} = 2 \end{cases}$}
\answerTableTall{$(4;\, 6)$}{$(-4;\, 6)$}{$(-4;\, -6)$}{$(4;\, -6)$}{$\left(-\dfrac{1}{4};\, 6\right)$}
\solbtn{5}
\taskBlock{Система лінійних нерівностей --- сесія 1.06, завдання №12}
\zadtask{Розв'яжіть систему нерівностей $\begin{cases} 9 + 3x < 18, \\ 2 - 3x > -10. \end{cases}$}
\answerTable{$(-\infty;\, 3)$}{$(3;\, 4)$}{$(4;\, +\infty)$}{$(3;\, +\infty)$}{$(-\infty;\, 4)$}
\solbtn{6}
\taskBlock{Система рівнянь --- сесія 2.06, завдання №13}
\zadtask{Розв'яжіть систему рівнянь $\begin{cases} xy - 2x = 3, \\ xy - 2x + y = 11. \end{cases}$}
\answerTableTall{$\left(\dfrac{1}{2};\, 8\right)$}{$(2;\, 8)$}{$\left(\dfrac{1}{2};\, -8\right)$}{$\left(-\dfrac{1}{2};\, 8\right)$}{$\left(8;\, \dfrac{1}{2}\right)$}
\solbtn{7}
\taskBlock{Система лінійних рівнянь --- сесія 4.06, завдання №4}
\zadtask{Розв'яжіть систему рівнянь $\begin{cases} x + 2y = 8, \\ x - 2y = -4. \end{cases}$}
\answerTable{$(6;\, 1)$}{$(3;\, 2)$}{$(-2;\, 5)$}{$(2;\, 3)$}{$(2;\, 4)$}
\solbtn{8}
\taskBlock{Система рівнянь --- сесія 5.06, завдання №4}
\zadtask{Розв'яжіть систему рівнянь $\begin{cases} \dfrac{8}{x} = -2, \\[0.2cm] x + \dfrac{y}{2} = 0. \end{cases}$}
\answerTable{$(-4;\, 8)$}{$(4;\, -8)$}{$(-4;\, -8)$}{$(4;\, 8)$}{$(-2;\, 4)$}
\solbtn{9}
\typeTitle{Квадратні нерівності}
\taskBlock{Квадратна нерівність --- сесія 23.05, завдання №11}
\zadtask{Знайдіть найбільший цілий розв'язок нерівності $(x - 4)(4 + x) \leq 8x - 7$. \nmtyear{2026}}
\answerTable{$0$}{$9$}{$-1$}{$10$}{$8$}
\solbtn{10}
\taskBlock{Квадратна нерівність --- сесія 5.06, завдання №11}
\zadtask{Розв'яжіть нерівність $(2x - 3)^2 < 25$.}
\answerTable{$(-1;\, 4)$}{$(-4;\, 1)$}{$(-1;\, +\infty)$}{$(-\infty;\, 4)$}{$(1;\, 4)$}
\solbtn{11}
\typeTitle{Ірраціональні рівняння}
\taskBlock{Ірраціональне рівняння --- сесія 30.05, завдання №14}
\zadtask{Розв'яжіть рівняння $\sqrt{12 - 2x} = 6$.}
\answerTable{$-12$}{$-3$}{$3$}{$12$}{$24$}
\solbtn{12}
\taskBlock{Ірраціональне рівняння --- сесія 1.06, завдання №14}
\zadtask{Розв'яжіть рівняння $\sqrt{\dfrac{25+2x}{9}} = \dfrac{1}{3}$.}
\answerTable{$-12$}{$-11$}{$-8$}{$0$}{$12$}
\solbtn{13}
\typeTitle{Показникові рівняння та нерівності}
\taskBlock{Система рівнянь (показникове + дробове) --- сесія 23.05, завдання №15}
\zadtask{Розв'яжіть систему рівнянь $\begin{cases} 5^x = \sqrt{5}, \\ \dfrac{y + 1}{2x + 3} = \dfrac{1 - 4x}{2}. \end{cases}$ Якщо $(x_0;\, y_0)$~--- розв'язок системи, то $x_0 + y_0 = $ \nmtyear{2026}}
\answerTable{$-2{,}5$}{$-1{,}5$}{$2{,}5$}{$-3{,}5$}{$1{,}5$}
\solbtn{14}
\taskBlock{Показникове рівняння --- сесія 1.06, завдання №1}
\zadtask{Позначте число, що є коренем рівняння $3 \cdot 3^x = 27$.}
\answerTable{$2$}{$3$}{$4$}{$9$}{$27$}
\solbtn{15}
\taskBlock{Показникове рівняння (проміжок для кореня) --- сесія 3.06, завдання №11}
\zadtask{Якому проміжку належить корінь рівняння $2^{x+1} - 2^x = \dfrac{1}{8}$?}
\answerTable{$(-\infty;\, -5)$}{$[-5;\, -3)$}{$[-3;\, 0)$}{$[0;\, 5)$}{$[5;\, +\infty)$}
\solbtn{16}
\taskBlock{Показникова нерівність --- сесія 4.06, завдання №15}
\zadtask{Розв'яжіть нерівність $3^{x^2 - 2x + 1} \leqslant 1$.}
\answerTable{$1$}{$(-\infty;\, -1]$}{$-1$}{$[1;\, +\infty)$}{$[-1;\, +\infty)$}
\solbtn{17}
\taskBlock{Показникове рівняння --- сесія 5.06, завдання №13}
\zadtask{Якому проміжку належить корінь рівняння $5^x \cdot 20^x = 1000$?}
\answerTable{$(-\infty;\, 0)$}{$[0;\, 1{,}5)$}{$[1{,}5;\, 2)$}{$[2;\, 4)$}{$[4;\, +\infty)$}
\solbtn{18}
\typeTitle{Логарифмічні рівняння та нерівності}
\taskBlock{Логарифмічна нерівність --- сесія 30.05, завдання №12}
\zadtask{Розв'яжіть нерівність $\log_5(4x - 7) < 2$.}
\answerTableTall
{$\left(\dfrac{7}{4};\, 8\right)$}
{$(8;\, +\infty)$}
{$\left(\dfrac{7}{4};\, +\infty\right)$}
{$(-\infty;\, 8)$}
{$\left(-\infty;\, \dfrac{7}{4}\right)$}
\solbtn{19}
\taskBlock{Логарифмічне рівняння --- сесія 2.06, завдання №10}
\zadtask{Якому проміжку належить корінь рівняння $\dfrac{\log_7(3x-4) - \log_7(x+5)}{2} = 0$?}
\answerTable{$[-1;\, 1)$}{$[1;\, 3)$}{$[3;\, 5)$}{$[5;\, 7)$}{$[7;\, +\infty)$}
\solbtn{20}
\taskBlock{Система рівнянь із логарифмом --- сесія 3.06, завдання №15}
\zadtask{Розв'яжіть систему рівнянь $\begin{cases} \log_2 x = 2, \\[0.15cm] \dfrac{y+10}{x+1} = \dfrac{x+2}{2}. \end{cases}$ Якщо $(x_0;\,y_0)$~--- розв'язок системи, то $x_0 + y_0 = {}$}
\answerTable{$18$}{$29$}{$-6$}{$-4$}{$9$}
\solbtn{21}
\taskBlock{Логарифмічне рівняння --- сесія 4.06, завдання №7}
\zadtask{Розв'яжіть рівняння $\log_9 (2x - 5) = \dfrac{1}{2}$.}
\answerTable{$-1$}{$-4$}{$43$}{$4$}{$16$}
\solbtn{22}
\typeTitle{Завдання з параметром (коротка відповідь)}
\taskBlock{Рівняння з параметром (відсутність коренів) --- сесія 23.05, завдання №22}
\zadtask{Знайдіть найбільше значення $a$, за якого рівняння $\dfrac{\sqrt{x^2 - 10x} - \sqrt{36 - 5x}}{x - 2\log_2 a} = 0$ не має коренів. \nmtyear{2026}}
\nmtAnswerBox
\solbtn{23}
\taskBlock{Показникове рівняння з параметром --- сесія 30.05, завдання №22}
\zadtask{Знайдіть \textit{кількість} усіх цілих значень $a$, за кожного з яких рівняння
$$9^x - (a+6) \cdot 3^x + 12a - 2a^2 = 0$$
має два різні корені.
\nmtAnswerBox}
\solbtn{24}
\taskBlock{Задача з параметром (два різні корені) --- сесія 1.06, завдання №22}
\zadtask{Знайдіть \textit{суму} всіх цілих значень параметра $a$ з проміжку $(0;\,7)$, за кожного з яких рівняння
$$\dfrac{x^2 + (a-3)x - a + 2}{2x - 3a + 6} = 0$$
має \textit{два різні} корені.}
\nmtAnswerBox
\solbtn{25}
\taskBlock{Рівняння з параметром --- сесія 2.06, завдання №22}
\zadtask{Визначте \textit{найменше} значення $a$, за якого рівняння
$$\dfrac{7^{\,3x - 2x^2} - \dfrac{1}{49}}{\sqrt{2x + 5a} - 4} = 0$$
має \textit{два різні} корені.
\nmtAnswerBox
}
\solbtn{26}
\taskBlock{Рівняння з параметром і модулем --- сесія 3.06, завдання №22}
\zadtask{Визначте \textit{суму} всіх значень $a$, за кожного з яких рівняння
$$\frac{|4x + 5a - 3| - 3}{(x+2)(x^2 + 9)} = 0$$
має \textit{єдиний} корінь.}
\nmtAnswerBox
\solbtn{27}
\taskBlock{Рівняння з параметром --- сесія 4.06, завдання №22}
\zadtask{Визначте \textit{найменше} ціле значення параметра $a$, за якого рівняння
\[ \frac{x + a + \sqrt{x+a} - 12}{2 + \lg^2(2x - a^2)} = 0 \]
має корені.
\nmtAnswerBox}
\solbtn{28}
\taskBlock{Рівняння з параметром --- сесія 5.06, завдання №22}
\zadtask{Знайдіть усі значення параметра $a$, за яких рівняння
$$\dfrac{\log_9(2x - 4a + 3) - 2}{\sqrt{a^2 - x^2}} = 0$$
має розв'язок.}
\nmtAnswerBox
\solbtn{29}
\chapterTitle{ПОЯСНЕННЯ ДО ЗАДАЧ}
\noindent\small Стислі пояснення до всіх оригінальних задач. Натискайте \fbox{\footnotesize $\leftarrow$ Задача №$N$}, щоб повернутися.\par\vspace{0.4cm}
\par\vspace{0.3cm}\noindent\hypertarget{sol:1}{}
\begin{tcolorbox}[colback=titleBg!40, colframe=mainGreen!55, boxrule=0.4pt, arc=2pt, left=6pt, right=6pt, top=4pt, bottom=4pt, breakable]
\noindent\small\textbf{Задача №1.}\quad Помноживши обидві частини на $7$: $x=9\cdot 7=63$. Відповідь: \textbf{Б}.
\par\vspace{0.1cm}\hfill\hyperlink{task:1}{\fbox{\footnotesize $\leftarrow$ Задача №1}}
\end{tcolorbox}
\par\vspace{0.3cm}\noindent\hypertarget{sol:2}{}
\begin{tcolorbox}[colback=titleBg!40, colframe=mainGreen!55, boxrule=0.4pt, arc=2pt, left=6pt, right=6pt, top=4pt, bottom=4pt, breakable]
\noindent\small\textbf{Задача №2.}\quad Нехай безе~--- $x$, тоді бісквітних~--- $5x$. Усього $x+5x=6x=30$, звідки $x=5$. Бісквітних: $5x=25$. Відповідь: \textbf{Б}.
\par\vspace{0.1cm}\hfill\hyperlink{task:2}{\fbox{\footnotesize $\leftarrow$ Задача №2}}
\end{tcolorbox}
\par\vspace{0.3cm}\noindent\hypertarget{sol:3}{}
\begin{tcolorbox}[colback=titleBg!40, colframe=mainGreen!55, boxrule=0.4pt, arc=2pt, left=6pt, right=6pt, top=4pt, bottom=4pt, breakable]
\noindent\small\textbf{Задача №3.}\quad Додаємо $3$ до всіх частин: $-8+3\leqslant x\leqslant5+3$, тобто $-5\leqslant x\leqslant8$, або $[-5;\,8]$. Відповідь: \textbf{Б}.
\par\vspace{0.1cm}\hfill\hyperlink{task:3}{\fbox{\footnotesize $\leftarrow$ Задача №3}}
\end{tcolorbox}
\par\vspace{0.3cm}\noindent\hypertarget{sol:4}{}
\begin{tcolorbox}[colback=titleBg!40, colframe=mainGreen!55, boxrule=0.4pt, arc=2pt, left=6pt, right=6pt, top=4pt, bottom=4pt, breakable]
\noindent\small\textbf{Задача №4.}\quad $2x - 1 \geqslant 0 \Rightarrow 2x \geqslant 1 \Rightarrow x \geqslant \dfrac12$, тобто $\left[\dfrac12;\, +\infty\right)$. Відповідь: \textbf{Д}.
\par\vspace{0.1cm}\hfill\hyperlink{task:4}{\fbox{\footnotesize $\leftarrow$ Задача №4}}
\end{tcolorbox}
\par\vspace{0.3cm}\noindent\hypertarget{sol:5}{}
\begin{tcolorbox}[colback=titleBg!40, colframe=mainGreen!55, boxrule=0.4pt, arc=2pt, left=6pt, right=6pt, top=4pt, bottom=4pt, breakable]
\noindent\small\textbf{Задача №5.}\quad З першого рівняння $\dfrac{2}{x}=-0{,}5\Rightarrow x=-4$; з другого $\dfrac{y}{3}=2\Rightarrow y=6$. Розв'язок $(-4;\,6)$. Відповідь: \textbf{Б}.
\par\vspace{0.1cm}\hfill\hyperlink{task:5}{\fbox{\footnotesize $\leftarrow$ Задача №5}}
\end{tcolorbox}
\par\vspace{0.3cm}\noindent\hypertarget{sol:6}{}
\begin{tcolorbox}[colback=titleBg!40, colframe=mainGreen!55, boxrule=0.4pt, arc=2pt, left=6pt, right=6pt, top=4pt, bottom=4pt, breakable]
\noindent\small\textbf{Задача №6.}\quad Перша: $3x<9 \Rightarrow x<3$. Друга: $-3x>-12 \Rightarrow x<4$. Перетин: $x<3$, тобто $(-\infty;\,3)$. Відповідь: \textbf{А}.
\par\vspace{0.1cm}\hfill\hyperlink{task:6}{\fbox{\footnotesize $\leftarrow$ Задача №6}}
\end{tcolorbox}
\par\vspace{0.3cm}\noindent\hypertarget{sol:7}{}
\begin{tcolorbox}[colback=titleBg!40, colframe=mainGreen!55, boxrule=0.4pt, arc=2pt, left=6pt, right=6pt, top=4pt, bottom=4pt, breakable]
\noindent\small\textbf{Задача №7.}\quad Віднявши перше рівняння від другого, дістаємо $y=11-3=8$. Підставляємо у перше: $8x-2x=3$, $6x=3$, $x=\dfrac{1}{2}$. Отже $\left(\dfrac{1}{2};\,8\right)$. Відповідь: \textbf{А}.
\par\vspace{0.1cm}\hfill\hyperlink{task:7}{\fbox{\footnotesize $\leftarrow$ Задача №7}}
\end{tcolorbox}
\par\vspace{0.3cm}\noindent\hypertarget{sol:8}{}
\begin{tcolorbox}[colback=titleBg!40, colframe=mainGreen!55, boxrule=0.4pt, arc=2pt, left=6pt, right=6pt, top=4pt, bottom=4pt, breakable]
\noindent\small\textbf{Задача №8.}\quad Додаємо рівняння: $2x=4$, $x=2$. Тоді $2+2y=8$, $2y=6$, $y=3$. Отже $(2;\,3)$. Відповідь: \textbf{Г}.
\par\vspace{0.1cm}\hfill\hyperlink{task:8}{\fbox{\footnotesize $\leftarrow$ Задача №8}}
\end{tcolorbox}
\par\vspace{0.3cm}\noindent\hypertarget{sol:9}{}
\begin{tcolorbox}[colback=titleBg!40, colframe=mainGreen!55, boxrule=0.4pt, arc=2pt, left=6pt, right=6pt, top=4pt, bottom=4pt, breakable]
\noindent\small\textbf{Задача №9.}\quad З першого рівняння $\dfrac{8}{x}=-2\Rightarrow x=-4$. Підставляємо в друге: $-4+\dfrac{y}{2}=0\Rightarrow y=8$. Розв'язок $(-4;\,8)$. Відповідь: \textbf{А}.
\par\vspace{0.1cm}\hfill\hyperlink{task:9}{\fbox{\footnotesize $\leftarrow$ Задача №9}}
\end{tcolorbox}
\par\vspace{0.3cm}\noindent\hypertarget{sol:10}{}
\begin{tcolorbox}[colback=titleBg!40, colframe=mainGreen!55, boxrule=0.4pt, arc=2pt, left=6pt, right=6pt, top=4pt, bottom=4pt, breakable]
\noindent\small\textbf{Задача №10.}\quad $(x-4)(x+4)\le 8x-7 \Rightarrow x^2-16\le 8x-7 \Rightarrow x^2-8x-9\le 0$. Дискримінант $D=64+36=100$, корені $x=\dfrac{8\pm 10}{2}$, тобто $-1$ і $9$. Отже $-1\le x\le 9$, найбільший цілий~--- $9$. Відповідь: \textbf{Б}.
\par\vspace{0.1cm}\hfill\hyperlink{task:10}{\fbox{\footnotesize $\leftarrow$ Задача №10}}
\end{tcolorbox}
\par\vspace{0.3cm}\noindent\hypertarget{sol:11}{}
\begin{tcolorbox}[colback=titleBg!40, colframe=mainGreen!55, boxrule=0.4pt, arc=2pt, left=6pt, right=6pt, top=4pt, bottom=4pt, breakable]
\noindent\small\textbf{Задача №11.}\quad $(2x-3)^2<25\iff |2x-3|<5\iff -5<2x-3<5\iff -2<2x<8\iff -1<x<4$. Відповідь: \textbf{А}.
\par\vspace{0.1cm}\hfill\hyperlink{task:11}{\fbox{\footnotesize $\leftarrow$ Задача №11}}
\end{tcolorbox}
\par\vspace{0.3cm}\noindent\hypertarget{sol:12}{}
\begin{tcolorbox}[colback=titleBg!40, colframe=mainGreen!55, boxrule=0.4pt, arc=2pt, left=6pt, right=6pt, top=4pt, bottom=4pt, breakable]
\noindent\small\textbf{Задача №12.}\quad Піднесемо до квадрата: $12-2x=36\Rightarrow -2x=24\Rightarrow x=-12$. Перевірка: $\sqrt{12-2\cdot(-12)}=\sqrt{36}=6$. Відповідь: \textbf{А} ($-12$).
\par\vspace{0.1cm}\hfill\hyperlink{task:12}{\fbox{\footnotesize $\leftarrow$ Задача №12}}
\end{tcolorbox}
\par\vspace{0.3cm}\noindent\hypertarget{sol:13}{}
\begin{tcolorbox}[colback=titleBg!40, colframe=mainGreen!55, boxrule=0.4pt, arc=2pt, left=6pt, right=6pt, top=4pt, bottom=4pt, breakable]
\noindent\small\textbf{Задача №13.}\quad Піднесемо до квадрата: $\dfrac{25+2x}{9}=\dfrac{1}{9}$, звідки $25+2x=1$, $2x=-24$, $x=-12$. Відповідь: \textbf{А}.
\par\vspace{0.1cm}\hfill\hyperlink{task:13}{\fbox{\footnotesize $\leftarrow$ Задача №13}}
\end{tcolorbox}
\par\vspace{0.3cm}\noindent\hypertarget{sol:14}{}
\begin{tcolorbox}[colback=titleBg!40, colframe=mainGreen!55, boxrule=0.4pt, arc=2pt, left=6pt, right=6pt, top=4pt, bottom=4pt, breakable]
\noindent\small\textbf{Задача №14.}\quad З першого рівняння $5^x=5^{1/2}$, тож $x_0=0{,}5$. Підставимо: $\dfrac{y+1}{2\cdot 0{,}5+3}=\dfrac{1-4\cdot 0{,}5}{2}$, тобто $\dfrac{y+1}{4}=-\dfrac{1}{2}$, звідки $y+1=-2$, $y_0=-3$. Тоді $x_0+y_0=0{,}5-3=-2{,}5$. Відповідь: \textbf{А}.
\par\vspace{0.1cm}\hfill\hyperlink{task:14}{\fbox{\footnotesize $\leftarrow$ Задача №14}}
\end{tcolorbox}
\par\vspace{0.3cm}\noindent\hypertarget{sol:15}{}
\begin{tcolorbox}[colback=titleBg!40, colframe=mainGreen!55, boxrule=0.4pt, arc=2pt, left=6pt, right=6pt, top=4pt, bottom=4pt, breakable]
\noindent\small\textbf{Задача №15.}\quad $3\cdot 3^x=27 \Rightarrow 3^{x+1}=3^3 \Rightarrow x+1=3 \Rightarrow x=2$. Відповідь: \textbf{А}.
\par\vspace{0.1cm}\hfill\hyperlink{task:15}{\fbox{\footnotesize $\leftarrow$ Задача №15}}
\end{tcolorbox}
\par\vspace{0.3cm}\noindent\hypertarget{sol:16}{}
\begin{tcolorbox}[colback=titleBg!40, colframe=mainGreen!55, boxrule=0.4pt, arc=2pt, left=6pt, right=6pt, top=4pt, bottom=4pt, breakable]
\noindent\small\textbf{Задача №16.}\quad $2^{x+1} - 2^x = 2^x(2-1) = 2^x = \dfrac18 = 2^{-3}$, звідки $x = -3 \in [-3;\, 0)$. Відповідь: \textbf{В}.
\par\vspace{0.1cm}\hfill\hyperlink{task:16}{\fbox{\footnotesize $\leftarrow$ Задача №16}}
\end{tcolorbox}
\par\vspace{0.3cm}\noindent\hypertarget{sol:17}{}
\begin{tcolorbox}[colback=titleBg!40, colframe=mainGreen!55, boxrule=0.4pt, arc=2pt, left=6pt, right=6pt, top=4pt, bottom=4pt, breakable]
\noindent\small\textbf{Задача №17.}\quad Оскільки $1=3^0$ і основа $3>1$, нерівність рівносильна $x^2-2x+1\leqslant0$, тобто $(x-1)^2\leqslant0$. Це виконується лише при $x=1$. Відповідь: \textbf{А}.
\par\vspace{0.1cm}\hfill\hyperlink{task:17}{\fbox{\footnotesize $\leftarrow$ Задача №17}}
\end{tcolorbox}
\par\vspace{0.3cm}\noindent\hypertarget{sol:18}{}
\begin{tcolorbox}[colback=titleBg!40, colframe=mainGreen!55, boxrule=0.4pt, arc=2pt, left=6pt, right=6pt, top=4pt, bottom=4pt, breakable]
\noindent\small\textbf{Задача №18.}\quad $5^x\cdot 20^x=(5\cdot 20)^x=100^x=10^{2x}$, а $1000=10^3$. Отже $10^{2x}=10^3\Rightarrow 2x=3\Rightarrow x=1{,}5$, що належить $[1{,}5;\,2)$. Відповідь: \textbf{В}.
\par\vspace{0.1cm}\hfill\hyperlink{task:18}{\fbox{\footnotesize $\leftarrow$ Задача №18}}
\end{tcolorbox}
\par\vspace{0.3cm}\noindent\hypertarget{sol:19}{}
\begin{tcolorbox}[colback=titleBg!40, colframe=mainGreen!55, boxrule=0.4pt, arc=2pt, left=6pt, right=6pt, top=4pt, bottom=4pt, breakable]
\noindent\small\textbf{Задача №19.}\quad ОДЗ: $4x-7>0$. Оскільки $2=\log_5 25$, нерівність $4x-7<25$. Разом: $0<4x-7<25\Rightarrow \dfrac74<x<8$. Відповідь: \textbf{А} $\left(\dfrac74;\,8\right)$.
\par\vspace{0.1cm}\hfill\hyperlink{task:19}{\fbox{\footnotesize $\leftarrow$ Задача №19}}
\end{tcolorbox}
\par\vspace{0.3cm}\noindent\hypertarget{sol:20}{}
\begin{tcolorbox}[colback=titleBg!40, colframe=mainGreen!55, boxrule=0.4pt, arc=2pt, left=6pt, right=6pt, top=4pt, bottom=4pt, breakable]
\noindent\small\textbf{Задача №20.}\quad Рівняння рівносильне $\log_7(3x-4)-\log_7(x+5)=0$, тобто $\log_7\dfrac{3x-4}{x+5}=0$. Звідси $3x-4=x+5$, $x=4{,}5\in[3;5)$. Відповідь: \textbf{В}.
\par\vspace{0.1cm}\hfill\hyperlink{task:20}{\fbox{\footnotesize $\leftarrow$ Задача №20}}
\end{tcolorbox}
\par\vspace{0.3cm}\noindent\hypertarget{sol:21}{}
\begin{tcolorbox}[colback=titleBg!40, colframe=mainGreen!55, boxrule=0.4pt, arc=2pt, left=6pt, right=6pt, top=4pt, bottom=4pt, breakable]
\noindent\small\textbf{Задача №21.}\quad Із $\log_2 x = 2$ маємо $x_0 = 4$. Тоді $\dfrac{y+10}{5} = \dfrac{6}{2} = 3$, звідки $y+10 = 15$, $y_0 = 5$. Отже $x_0 + y_0 = 9$. Відповідь: \textbf{Д}.
\par\vspace{0.1cm}\hfill\hyperlink{task:21}{\fbox{\footnotesize $\leftarrow$ Задача №21}}
\end{tcolorbox}
\par\vspace{0.3cm}\noindent\hypertarget{sol:22}{}
\begin{tcolorbox}[colback=titleBg!40, colframe=mainGreen!55, boxrule=0.4pt, arc=2pt, left=6pt, right=6pt, top=4pt, bottom=4pt, breakable]
\noindent\small\textbf{Задача №22.}\quad За означенням логарифма $2x-5=9^{1/2}=3$, звідки $2x=8$, $x=4$ (ОДЗ $2x-5>0$ виконується). Відповідь: \textbf{Г}.
\par\vspace{0.1cm}\hfill\hyperlink{task:22}{\fbox{\footnotesize $\leftarrow$ Задача №22}}
\end{tcolorbox}
\par\vspace{0.3cm}\noindent\hypertarget{sol:23}{}
\begin{tcolorbox}[colback=titleBg!40, colframe=mainGreen!55, boxrule=0.4pt, arc=2pt, left=6pt, right=6pt, top=4pt, bottom=4pt, breakable]
\noindent\small\textbf{Задача №23.}\quad Чисельник нуль при $\sqrt{x^2-10x}=\sqrt{36-5x}$. ОДЗ: $x^2-10x\ge 0$ і $36-5x\ge 0$, тобто $x\le 0$ (бо $x\le 7{,}2$). Піднесемо до квадрата: $x^2-10x=36-5x \Rightarrow x^2-5x-36=0$; $D=25+144=169$, корені $x=\dfrac{5\pm 13}{2}=9$ або $-4$. В ОДЗ лише $x=-4$. Щоб рівняння не мало коренів, цей корінь має обнуляти знаменник: $-4=2\log_2 a$, тобто $\log_2 a=-2$, $a=0{,}25$. Відповідь: \textbf{0,25}.
\par\vspace{0.1cm}\hfill\hyperlink{task:23}{\fbox{\footnotesize $\leftarrow$ Задача №23}}
\end{tcolorbox}
\par\vspace{0.3cm}\noindent\hypertarget{sol:24}{}
\begin{tcolorbox}[colback=titleBg!40, colframe=mainGreen!55, boxrule=0.4pt, arc=2pt, left=6pt, right=6pt, top=4pt, bottom=4pt, breakable]
\noindent\small\textbf{Задача №24.}\quad Заміна $t=3^x>0$: $t^2-(a+6)t+(12a-2a^2)=0$. За теоремою Вієта корені $t_1=2a$, $t_2=6-a$ ($t_1t_2=2a(6-a)=12a-2a^2$, $t_1+t_2=a+6$). Для двох різних коренів $x$ потрібні два різні додатні $t$: $2a>0$, $6-a>0$, $2a\neq6-a$, тобто $a\in(0;6)$, $a\neq2$. Цілі: $1,3,4,5$. Відповідь: $4$.
\par\vspace{0.1cm}\hfill\hyperlink{task:24}{\fbox{\footnotesize $\leftarrow$ Задача №24}}
\end{tcolorbox}
\par\vspace{0.3cm}\noindent\hypertarget{sol:25}{}
\begin{tcolorbox}[colback=titleBg!40, colframe=mainGreen!55, boxrule=0.4pt, arc=2pt, left=6pt, right=6pt, top=4pt, bottom=4pt, breakable]
\noindent\small\textbf{Задача №25.}\quad Чисельник: $x^2+(a-3)x-(a-2)=0$. За теоремою Вієта корені мають суму $3-a$ і добуток $2-a$; підбором (Вієта) це $x_1=1$ і $x_2=2-a$. Для двох різних коренів потрібно $x_1\ne x_2$, тобто $a\ne 1$, і жоден корінь не має обертати знаменник $2x-3a+6$ на нуль; знаменник занулюється коренем $x_2$ при $a=2$. Цілі $a\in(0;7)$: $1,2,3,4,5,6$; вилучаємо $a=1$ і $a=2$. Сума $3+4+5+6=18$. Відповідь: \textbf{$18$}.
\par\vspace{0.1cm}\hfill\hyperlink{task:25}{\fbox{\footnotesize $\leftarrow$ Задача №25}}
\end{tcolorbox}
\par\vspace{0.3cm}\noindent\hypertarget{sol:26}{}
\begin{tcolorbox}[colback=titleBg!40, colframe=mainGreen!55, boxrule=0.4pt, arc=2pt, left=6pt, right=6pt, top=4pt, bottom=4pt, breakable]
\noindent\small\textbf{Задача №26.}\quad Чисельник $=0$: $7^{3x-2x^2}=7^{-2}$, тобто $3x-2x^2=-2$, або $2x^2-3x-2=0$. За теоремою Вієта корені $x_1=2$, $x_2=-\dfrac{1}{2}$ (бо $x_1+x_2=\dfrac{3}{2}$, $x_1x_2=-1$). Щоб обидва корені були допустимі, потрібно $2x+5a\geqslant0$ у меншому корені $x=-\dfrac{1}{2}$: $-1+5a\geqslant0$, звідки $a\geqslant0{,}2$. Найменше $a=0{,}2$.
\par\vspace{0.1cm}\hfill\hyperlink{task:26}{\fbox{\footnotesize $\leftarrow$ Задача №26}}
\end{tcolorbox}
\par\vspace{0.3cm}\noindent\hypertarget{sol:27}{}
\begin{tcolorbox}[colback=titleBg!40, colframe=mainGreen!55, boxrule=0.4pt, arc=2pt, left=6pt, right=6pt, top=4pt, bottom=4pt, breakable]
\noindent\small\textbf{Задача №27.}\quad Оскільки $x^2+9>0$, ОДЗ: $x\ne-2$. Рівняння $|4x+5a-3| = 3$ дає $4x+5a-3 = \pm3$, тобто два корені $x_1 = \dfrac{6-5a}{4}$ і $x_2 = \dfrac{-5a}{4}$ ($x_1\ne x_2$). Єдиний корінь буде, коли один із них дорівнює $-2$: $x_1=-2 \Rightarrow a=2{,}8$; $x_2=-2 \Rightarrow a=1{,}6$. Сума: $2{,}8+1{,}6 = 4{,}4$. Відповідь: \textbf{$4{,}4$}.
\par\vspace{0.1cm}\hfill\hyperlink{task:27}{\fbox{\footnotesize $\leftarrow$ Задача №27}}
\end{tcolorbox}
\par\vspace{0.3cm}\noindent\hypertarget{sol:28}{}
\begin{tcolorbox}[colback=titleBg!40, colframe=mainGreen!55, boxrule=0.4pt, arc=2pt, left=6pt, right=6pt, top=4pt, bottom=4pt, breakable]
\noindent\small\textbf{Задача №28.}\quad Знаменник завжди додатний, тож потрібен нуль чисельника. Нехай $t=\sqrt{x+a}\geqslant0$: $t^2+t-12=0$. За теоремою Вієта $t_1=3$, $t_2=-4$ ($t_1+t_2=-1$, $t_1t_2=-12$); підходить $t=3$, отже $x+a=9$, $x=9-a$. Має існувати $\lg(2x-a^2)$, тобто $2x-a^2>0$: $18-2a-a^2>0$, або $a^2+2a-18<0$. Корені (через дискримінант $D=4+72=76$): $a=-1\pm\sqrt{19}$, тож $-1-\sqrt{19}<a<-1+\sqrt{19}$, тобто приблизно $-5{,}36<a<3{,}36$. Найменше ціле $a=-5$. Відповідь: \textbf{$-5$}.
\par\vspace{0.1cm}\hfill\hyperlink{task:28}{\fbox{\footnotesize $\leftarrow$ Задача №28}}
\end{tcolorbox}
\par\vspace{0.3cm}\noindent\hypertarget{sol:29}{}
\begin{tcolorbox}[colback=titleBg!40, colframe=mainGreen!55, boxrule=0.4pt, arc=2pt, left=6pt, right=6pt, top=4pt, bottom=4pt, breakable]
\noindent\small\textbf{Задача №29.}\quad Чисельник $=0$: $\log_9(2x-4a+3)=2\Rightarrow 2x-4a+3=81\Rightarrow x=2a+39$. Знаменник вимагає $a^2-x^2>0$, тобто $-|a|<x<|a|$, що можливо лише при $a<0$. Умова $|2a+39|<-a$ дає систему $a<2a+39<-a$, звідки $-39<a$ і $a<-13$. Відповідь: \textbf{$a\in(-39;\,-13)$}.
\par\vspace{0.1cm}\hfill\hyperlink{task:29}{\fbox{\footnotesize $\leftarrow$ Задача №29}}
\end{tcolorbox}
\chapterTitle{ВІДПОВІДІ}
\begin{multicols}{3}\noindent
\textbf{1.}~Б\\
\textbf{2.}~Б\\
\textbf{3.}~Б\\
\textbf{4.}~Д\\
\textbf{5.}~Б\\
\textbf{6.}~А\\
\textbf{7.}~А\\
\textbf{8.}~Г\\
\textbf{9.}~А\\
\textbf{10.}~Б\\
\textbf{11.}~А\\
\textbf{12.}~А\\
\textbf{13.}~А\\
\textbf{14.}~А\\
\textbf{15.}~А\\
\textbf{16.}~В\\
\textbf{17.}~А\\
\textbf{18.}~В\\
\textbf{19.}~А\\
\textbf{20.}~В\\
\textbf{21.}~Д\\
\textbf{22.}~Г\\
\textbf{23.}~$0{,}25$\\
\textbf{24.}~$4$\\
\textbf{25.}~$18$\\
\textbf{26.}~$0{,}2$\\
\textbf{27.}~$4{,}4$\\
\textbf{28.}~$-5$\\
\textbf{29.}~$a\in(-39;\,-13)$\\
\end{multicols}
\end{document}
